\documentclass[10pt]{article}
\usepackage[T1]{fontenc}
\usepackage[utf8]{inputenc}
\usepackage[french]{babel}
\usepackage{amsmath,amssymb,mathrsfs,mathpazo}
\usepackage{geometry}
\geometry{a4paper,margin=2cm}
\usepackage{graphicx,array,tabularx,tabularray,colortbl,multirow,multicol,booktabs,enumitem,pifont,ulem,fancybox,diagbox,cancel}
\usepackage{pstricks,pstricks-add,pst-plot,pst-tree,pst-node,pst-text,pst-func,pst-math,pst-bspline,pst-3dplot,pst-eucl,pst-solides3d}
\let\pscalculate\undefined
\usepackage{pst-bezier}
\usepackage{tikz,pgfplots,tkz-tab}
\usetikzlibrary{arrows,arrows.meta,patterns,calc,positioning,trees,decorations.markings,shapes.geometric,shapes.misc,shapes.arrows,intersections,decorations.pathmorphing,decorations.pathreplacing,backgrounds,shadows,babel}
\pgfplotsset{compat=1.18}
\usepgfplotslibrary{fillbetween}
\usepackage[np]{numprint}
\usepackage{listings,makecell,pas-tableur,twoopt}
\let\barre\undefined
\usepackage[pstricks]{tablvar}
\definecolor{darkgreen}{RGB}{0,100,0}
\providecommand{\red}{\color{red}}
\providecommand{\blue}{\color{blue}}
\providecommand{\green}{\color{green}}
\providecommand{\black}{\color{black}}
\providecommand{\white}{\color{white}}
\providecommand{\gray}{\color{gray}}
\providecommand{\cyan}{\color{cyan}}
\providecommand{\magenta}{\color{magenta}}
\providecommand{\yellow}{\color{yellow}}
\providecommand{\violet}{\color{violet}}
\providecommand{\orange}{\color{orange}}
\providecommand{\euro}{\texteuro}
\providecommand{\EUR}{\texteuro}
\providecommand{\ieme}{\textsuperscript{e}}
\providecommand{\degres}{\ensuremath{^\circ}}
\providecommand{\degree}{\ensuremath{^\circ}}
\providecommand{\vv}[1]{\overrightarrow{#1}}
\providecommand{\Coord}[1]{#1}
\providecommand{\decosix}{}
\providecommand{\starredbullet}{\ensuremath{\bullet}}
\pagestyle{empty}
\setlength{\parindent}{0pt}
\usepackage{adjustbox}\DeclareUnicodeCharacter{2236}{:}\DeclareUnicodeCharacter{2190}{\ensuremath{\leftarrow}}
\providecommand{\R}{\mathbb{R}}
\providecommand{\N}{\mathbb{N}}
\providecommand{\Z}{\mathbb{Z}}
\providecommand{\Q}{\mathbb{Q}}
\providecommand{\D}{\mathbb{D}}
\providecommand{\C}{\mathbb{C}}
\providecommand{\e}{{\mathrm{e}}}
\providecommand{\vect}[1]{\overrightarrow{#1}}
\providecommand{\vectt}[1]{\overrightarrow{\mathrm{#1}}}
\providecommand{\vvt}[1]{\overrightarrow{\mathrm{#1}}}
\providecommand{\np}[2][]{{#2~#1}}
\providecommand{\vv}[1]{{\overrightarrow{#1}}}
\renewcommand{\d}{{\mathrm{d}}}
\providecommand{\degres}{{^\circ}}
\providecommand{\degre}{{^\circ}}
\providecommand{\blue}{\color{blue}}
\providecommand{\red}{\color{red}}
\providecommand{\enskip}{\quad}
\providecommand{\llbracket}{\lbrack\!\lbrack}
\providecommand{\rrbracket}{\rbrack\!\rbrack}
\providecommand{\euro}{{\text{€}}}
\providecommand{\no}{{\text{n°}}}
\providecommand{\interval}[2]{{[#1;#2]}}
\providecommand{\intervFO}[2]{{[#1;#2[}}
\providecommand{\dotfill}{\ldots}

\providecommand{\ds}{\displaystyle}
\providecommand{\textding}[1]{#1}
\providecommand{\Oij}{\(\mathrm{O};\vec{\imath},\vec{\jmath}\)}
\providecommand{\Oijk}{\(\mathrm{O};\vec{\imath},\vec{\jmath},\vec{k}\)}
\providecommand{\Ouv}{\(\mathrm{O};\vec{u},\vec{v}\)}
\providecommand{\intervFF}[2]{\ensuremath{[#1;#2]}}
\renewcommand{\euro}{\ifmmode\text{\texteuro}\else\texteuro\fi}
\DeclareUnicodeCharacter{00B0}{\ifmmode^{\circ}\else\textdegree\fi}
\DeclareUnicodeCharacter{20AC}{\ifmmode\text{\texteuro}\else\texteuro\fi}
% Énoncés : Éducation nationale. Transcription et corrigés : APMEP, Denis Vergès et contributeurs. Découpage et classement : CoopMaths, Sébastien Lozano, Éric Elter et contributeurs. Adaptation de la présentation : Maths Explorer.
% CC BY-SA : https://coopmaths.fr/www/about/
% https://coopmaths.fr/alea/static/e3c/2020/tex/e3c_2020_00_sujet24_1_cor.tex
\begin{document}
\noindent\textit{APMEP / CoopMaths. CC BY-SA. Présentation adaptée par Maths Explorer.}\par\medskip
\subsection*{Question 1}
	
	Un vecteur directeur de la droite est 
	$\vec u\begin{pmatrix} -5 \\ 4 \end{pmatrix}$, donc un vecteur normal est par exemple $\vec{n}  \begin{pmatrix} 4 \\ 5 \end{pmatrix}$.
	
	\subsection*{Question 2}
	
	Le point $H(3 ; 4)$ appartient à la droite $(d)$ si et seulement si :
	\[
	4 \times 3 + 5 \times 4 - 32 = 0
	\]
	ce qui est vrai.
	
	De plus, le vecteur $\overrightarrow{AH} \begin{pmatrix} -4 \\ -5 \end{pmatrix}$ est bien orthogonal au vecteur directeur de la droite d'équation $4x + 5y - 32 = 0$, soit $\vec d \begin{pmatrix} -5 \\ 4 \end{pmatrix}$.
	
	Ainsi, le vecteur $\overrightarrow{AH} \begin{pmatrix} -4 \\ -5 \end{pmatrix}$ est un vecteur normal à la droite.
	
	\subsection*{Question 3}
	
	Le point $M(x ; y)$ appartient au cercle $\mathcal{C}(A, R = 2)$ si et seulement si :
	\[
	AM^2 = 2^2 \quad \text{soit} \quad (x - (-1))^2 + (y - 3)^2 = 4,
	\]
	ce qui équivaut à :
	\[
	(x + 1)^2 + (y - 3)^2 = 4.
	\]
	
	\subsection*{Question 4}
	
	L'équation de la parabole est donnée par :
	\[
	y = 3x^2 - 9x + 5.
	\]
	On peut la réécrire sous forme canonique :
	\[
	y = 3 \left( x^2 - 3x \right) + 5 = 3 \left[ \left( x - \dfrac{3}{2} \right)^2 - \dfrac{9}{4} \right] + 5 = 3 \left( x - \dfrac{3}{2} \right)^2 - \dfrac{27}{4} + 5 = 3 \left( x - \dfrac{3}{2} \right)^2 - \dfrac{7}{4}.
	\]
	Cette forme montre que le minimum de la courbe est atteint pour $x = \dfrac{3}{2}$ et que le minimum vaut $f\left( \dfrac{3}{2} \right) = -\dfrac{7}{4}$.
	
	Ainsi, le sommet de la parabole est $S\left( \dfrac{3}{2} ; -\dfrac{7}{4} \right)$ et l'axe de symétrie a pour équation $x = \dfrac{3}{2}$.
	
	\subsection*{Question 5}
	
	Le trinôme $-3x^2 + 9x - 5$ a un discriminant :
	\[
	\Delta = 81 - 4 \times (-3) \times (-5) = 81 - 60 = 21.
	\]
	Comme $\Delta > 0$, le trinôme a deux racines :
	\[
	x_1 = \dfrac{-9 + \sqrt{21}}{-6} \approx 0,74 \quad \text{et} \quad x_2 = \dfrac{-9 - \sqrt{21}}{-6} \approx 2,26.
	\]
	Puisque $a = -3 < 0$, le trinôme est négatif sur les intervalles en dehors des racines, donc il est négatif sur l'intervalle $[x_1 ; x_2]$.

\end{document}
